\documentclass{article}
\usepackage[T1]{fontenc}
\usepackage{lmodern}
\usepackage{graphicx}
\usepackage{amsmath,amssymb}
\usepackage[a4paper,margin=2cm]{geometry}
\usepackage[absolute,overlay]{textpos}
\setlength{\TPHorizModule}{1mm}
\setlength{\TPVertModule}{1mm}
\usepackage[indonesian]{babel}
\usepackage{hyperref}
\setlength{\emergencystretch}{2em}
% unit: Halaman 31

\begin{document}

% === Halaman 31 (python/native) ===

S-box

Elemen-elemen di dalam S-box terlihat seperti diisi secara acak, namun 

sebenarnya ia dihasilkan dari proses perhitungan sebagai berikut:

1. Inisialisasi S-box dengan nilai yang menaik dari baris ke baris. Baris 

ke-0 diisi dengan nilai 00, 01, 02, …, 0F. Baris ke-1 diisi dengan nilai 

10, 11, 12, …, 1F, dan seterusnya untuk baris-baris lain.

2. Untuk setiap nilai pada baris y kolom x, tentukan balikannya dalam 

GF(28). Nilai 00 dipetakan ke dirinya sendiri. Penjelasan tentang 

GF(28) lihat pada lampiran.

3. Hasil dari langkah 2 dikonversi ke vektor kolom bit (b0, b1, …, b8)T.

\end{document}
